\begin{proposition}[2.4.9]
\label{IV.2.4.9-fr}
Sea $Y$ un preesquema \emph{discret}o. Entonces todo morfismo $f:X\to Y$ es
universalmente abierto.
\end{proposition}

Como la cuestión es local en $Y$ (2.4.3, (vii)), es posible limitarse al caso
en que el espacio subyacente a $Y$ se reduce a un punto; reemplazando $f$ por
$f_{\mathrm{red}}$ (2.4.3, (vi)), se puede suponer además que $Y$ es el
espectro de un \emph{cuerpo} $k$; por otra parte, para todo cambio de base
$Y'\to Y$, los abiertos de $X'=X\times_Y Y'$ que son imágenes inversas de los
abiertos afines de $X$ recubren $X'$, por lo que se puede suponer que
$X=\operatorname{Spec}(B)$ es afín, siendo $B$ una \emph{$k$-álgebra}. Se trata,
pues, de demostrar que, para toda $k$-álgebra $A'$, si se pone
$B'=B\otimes_k A'$, la imagen por
$f':\operatorname{Spec}(B')\to\operatorname{Spec}(A')$ de todo abierto de la
forma $U=D(t)$ ($t\in B'$) es abierta en $\operatorname{Spec}(A')$. Ahora bien,
$B$ es el límite inductivo de la familia filtrante creciente $(B_\alpha)$ de
sus sub-$k$-álgebras \emph{de tipo finito}; por tanto (como el funtor
$\varinjlim$ conmuta con el producto tensorial), $B'$ es el límite inductivo
de las $A'$-álgebras $B_\alpha\otimes_k A'=B'_\alpha$; existe $\alpha$ tal que
$t$ sea la imagen en $B'$ de un elemento $t_\alpha\in B'_\alpha$, por lo que
$D(t)$ es la imagen inversa por el morfismo canónico
$u_\alpha:\operatorname{Spec}(B')\to\operatorname{Spec}(B'_\alpha)$ del
abierto $U_\alpha=D(t_\alpha)$ (\textbf I, 1.2.2.2). Pero, como $k$ es un cuerpo,
$B_\alpha$ es una $k$-álgebra de presentación finita y un $k$-módulo plano,
por lo que $B'_\alpha$ es una $A'$-álgebra de presentación finita y un
$A'$-módulo plano; se deduce, pues, de (2.4.6) que la imagen de $U_\alpha$ por
$f'_\alpha:\operatorname{Spec}(B'_\alpha)\to\operatorname{Spec}(A')$ es
\emph{abierta} en $\operatorname{Spec}(A')$. Todo se reduce, por tanto, a ver
que $f'(U)=f'_\alpha(U_\alpha)$; está claro que
$f'(U)\subset f'_\alpha(U_\alpha)$ y solo queda comprobar que, para todo punto
$y'\in f'_\alpha(U_\alpha)$, la intersección
$V=U\cap f'^{-1}(y')$ no es vacía. Ahora bien, se tiene
$V=u_\alpha^{-1}(V_\alpha)$, donde
$V_\alpha=U_\alpha\cap f_\alpha'^{-1}(y')$; en otras palabras, $V_\alpha$ es
un abierto no vacío (por definición de $y'$) del preesquema
$\operatorname{Spec}(B'_\alpha\otimes_{A'}\boldsymbol{k}(y'))
=\operatorname{Spec}(B_\alpha\otimes_k \boldsymbol{k}(y'))$
y $V$ es su imagen inversa por el morfismo
$v:\operatorname{Spec}(B\otimes_k \boldsymbol{k}(y'))
\longrightarrow\operatorname{Spec}(B_\alpha\otimes_k \boldsymbol{k}(y'))$.
Como $B_\alpha$ es una subálgebra de $B$ y $k$ es un cuerpo, el homomorfismo
$B_\alpha\otimes_k \boldsymbol{k}(y')\to B\otimes_k \boldsymbol{k}(y')$
\oldpage[IV-2]{22}
es \emph{inyectivo} por planitud, por lo que el morfismo $v$ es \emph{dominante}
(\textbf I, 1.2.7), lo que termina la demostración.

\begin{corollary}[2.4.10]
\label{IV.2.4.10-fr}
Sean $k$ un cuerpo y $X$, $Y$ dos $k$-preesquemas; entonces el morfismo de
proyección $X\times_k Y\to X$ es universalmente abierto. En particular,
para toda extensión $K$ de $k$, el morfismo de proyección $X_{(K)}\to X$ es
universalmente abierto.
\end{corollary}

Basta aplicar (2.4.9) al morfismo estructural
$Y\to\operatorname{Spec}(k)$.

\begin{remark}[2.4.11]
\label{IV.2.4.11-fr}
Si $f:X\to Y$ es un morfismo \emph{abierto}, se sabe que, para \emph{toda} parte
$E$ de $Y$, se tiene
$f^{-1}(\overline E)=\overline{f^{-1}(E)}$
(Bourbaki, \emph{Top. gén.}, cap.~I, 3.ª ed., \S~5,
núm.~4, prop.~7). Esta observación se aplica, por ejemplo,
cuando $f$ es un morfismo \emph{plano localmente de presentación finita}
(2.4.6), o un morfismo de proyección $X\times_k Y\to X$ donde $X$, $Y$ son
preesquemas sobre un cuerpo $k$ (2.4.10), y generaliza entonces (2.3.10).
\end{remark}
